\section{Related Work and Motivation}
\label{sec:motivation}

\subsection{Video Analytics System and VLMs}
\label{subsec:vas_vlm}



The field of video analytics has seen significant advancements in recent years~\cite{khani2023recl, bhardwaj2022ekya, zhang2024vulcan, padmanabhan2023gemel, lu2022turbo, sivaraman2024gemino, zhang2021sensei, jain20sec}. Leveraging emerging deep learning models, modern video analytics systems can extract insightful information, such as object locations or counts, from video streams processed on edge or cloud servers. 



Existing video analytics systems predominantly support \emph{closed-end} analytics (L1 to L3 systems as mentioned in \S\ref{sec:intro}), often relying on shallow neural networks to extract predefined, task-specific, and constrained information. For instance, Remix~\cite{remix_mobicom_21} leverages fine-tuned EfficientDet~\cite{tan2020efficientdet} to generate bounding boxes for pedestrians. Consequently, the flexibility and adaptability of video analytics systems are fundamentally limited by the inherent constraints of the specific models they employ.

Recently, VLMs like GPT-4o~\cite{achiam2023gpt}, Gemini~\cite{team2023gemini,team2024gemini}, QwenVL~\cite{wang2024qwen2,bai2025qwen2_5} and Phi~\cite{abdin2024phi3,abdin2024phi4,abouelenin2025phi4mini}, have demonstrated their transformative potential in video analytics tasks. By leveraging the extensive world knowledge embedded in large language models (LLMs), VLMs not only achieve generalized visual grounding but also exhibit advanced video comprehension capabilities, such as zero-shot temporal and spatial reasoning, contextual retrieval, and semantic understanding. More importantly, VLMs enable natural language interaction, allowing users to dynamically query video content. This makes them particularly well-suited for addressing diverse and unstructured \emph{open-ended} analytics in real-world scenarios, towards L4 video analytics systems.







However, adopting VLMs in video analytics systems is far from straightforward. Existing L1 to L3 video analytics systems, which rely on traditional DNNs, typically process each video frame independently. In contrast, VLMs require related frames to be processed collectively to infer causal relationships and temporal dependencies across frames. This shift introduces significant complexity, as existing VLMs are generally capable of handling only minute-level or sub-hour-level videos due to the limited context window inherent in language models.

In real-world video analytics scenarios, the scale of videos to be analyzed is often vastly larger—spanning hundreds of hours or more (\eg monitoring wildlife behavior over an entire month, as illustrated in Fig.~\ref{fig:evg_kg}). This creates a fundamental gap between the capabilities of VLMs and the demands of video analytics systems, as the limited context window length of VLMs directly restricts their ability to process videos of such extensive durations effectively.


\begin{table}[t]
\centering
\resizebox{\columnwidth}{!}{%
\begin{tabular}{cc|cc|cc} 
\toprule[2pt]
\multicolumn{2}{c}{\textbf{Short (1.4 minutes)}} & \multicolumn{2}{c}{\textbf{Medium (9.7 minutes)}} & \multicolumn{2}{c}{\textbf{Long (39.7 minutes)}}  \\ 
\cmidrule{1-6}
Total & Needed & Total & Needed   & Total & Needed \\ 
2144.8 & 12.1 (0.5\%)      & 13924.1     & 68.1 (0.4\%)                   & 66847.1     & 82.3 (0.1\%)                \\
\bottomrule[2pt]
\end{tabular}%
}
\caption{\revised{Only a small portion of the frames are necessary to answer each particular question across the short, medium, and long video subsets of VideoMME~\cite{fu2024videomme} benchmark using Qwen2-VL.}}
\label{tab:few_needed_frames}
\end{table}

\subsection{Long Video Understandings}




Recent research also increasingly focuses on enabling long video understandings~\cite{shen2024longvu, wang2025adaretake}. \revised{Current autoregressive language models inherently have the constrained context window length, therefore efforts have been directed toward reducing the redundancy in video inputs to facilitate the processing of extended video durations.  For instance, LongVU~\cite{shen2024longvu} and AdaRETAKE~\cite{wang2025adaretake} introduce dynamic compression mechanisms that prioritize video content based on its relevance, selectively retaining frames or regions most pertinent to downstream language tasks. Similarly, NVILA~\cite{liu2024nvila} addresses the efficiency-accuracy trade-off by optimizing sampling strategies and resolution to fit within limited token budgets.} 


While these approaches have succeeded in increasing the number of frames that models can process and mitigating the constraints of context windows to some extent, they fall short of achieving a fundamental breakthrough. \revised{Given that SOTA LLMs typically support context lengths ranging from 128K to 1M tokens, most existing approaches are thereby restricted to processing video segments of at most one hour in duration~\cite{qian2024streaming}, which is inadequate for the requirements of video analytics. Furthermore, as video length increases, the associated inference cost rises proportionally, thereby compounding the scalability challenges inherent to these systems.}











\subsection{Retrieval Augmented Generation}

\begin{figure*}[t]
   \centering
   \includegraphics[width=.85\linewidth]{Figures/overview/overview-crop.pdf}
   \caption{The system overview of \sysname.}
   \label{fig:overview}
\end{figure*}








While the videos to be analyzed may span extensive durations, the frames necessary to respond to a specific query are often limited. To validate this observation, we conducted an experiment on VideoMME~\cite{ma2024drvideo} using Qwen2-VL~\cite{wang2024qwen2}. Specifically, we first identified all questions for which Qwen2-VL produced correct answers by uniformly sampling frames from the videos at a rate of 1 FPS. For these questions, we then determined the minimal set of frames required for the VLM to generate the correct answer by iteratively reducing the number of input frames using a binary search strategy.~\footnote{For example, we initially uniformly sample 100 frames from the video. If the VLM can generate the correct answer based on these frames, we then attempt to reduce the frame set to 50 frames. If 50 frames are still sufficient to produce the correct answer, we further reduce the set to 25 frames. Conversely, if 50 frames are insufficient, we increase the set to 75 frames, iteratively refining the frame selection using binary search strategy.} \revised{The results indicate that, owing to significant temporal redundancy in video data, the frames required to answer a specific query represent only a small fraction of the total frames, as shown in Table~\ref{tab:few_needed_frames}.}


Based on this observation, an intuitive approach would be to first retrieve the relevant frames corresponding to a specific query and then generate the final answer based on these frames, a method commonly referred to as Retrieval-Augmented Generation (RAG). To retrieve potentially relevant frames, a straightforward strategy involves \emph{vectorized retrieval}, where each frame of the video is embedded using a vision-language model, such as CLIP~\cite{radford2021CLIP}. At query time, the embedding of the query is used to retrieve relevant frames by comparing the similarity between the query embedding and the vectorized frames. 

However, the vectorized retrieval method's limitations stem from the detailed information contained in the query. Notably, such an approach struggles to handle query-focused summaries~\cite{edge2024graphrag} (\eg "What happened in the last few hours?") or multi-hop queries~\cite{zhao2024retrieval} (\eg "What did the man do after he opened the fridge?"), as the retrieved frames often fail to capture key contexts that are not explicitly mentioned in the query descriptions.


To enable effective retrieval, recent research has explored two prominent approaches: \emph{video structuring} and \emph{iterative retrieval}. For example, Video-RAG~\cite{luo2024video-rag} structures videos by utilizing various tools to extract information such as automatic speech recognition (ASR), optical character recognition (OCR) results, and object detection outputs. It then applies RAG techniques to the structured information. However, this method is inherently constrained by the tools employed for video structuring. It is often impractical to predict in advance what types of information need to be extracted and what corresponding tools should be utilized, limiting its adaptability to diverse and dynamic video analytics scenarios.

Alternatively, researchers have proposed obtaining relevant frames through multiple iterative retrieval processes~\cite{wang2024videoagent, wang2024videotree, ma2024drvideo, yang2024vca}. For instance, VideoAgent~\cite{wang2024videoagent} typically begins with a coarse-grained sampling of video segments to establish an initial high-level understanding. Based on this, the VLM is prompted to decide which finer-grained segments to retrieve and analyze in subsequent iterations.  However, these approaches face significant challenges when applied to video analytics scenarios involving extremely long videos. On one hand, the initial coarse-grained sampling may become insufficient as video length increases, potentially missing critical information. On the other hand, the iterative retrieval and analysis process becomes increasingly computationally expensive as video duration grows, making it impractical for large-scale video analytics tasks.


Recent studies have advanced RAG techniques~\cite{edge2024graphrag, lazygraphrag2024, guo2024lightrag, fan2025minirag} by incorporating knowledge graph construction to enhance the retrieval process. However, these works primarily focus on text-only RAG problems, and adapting such approaches to video analytics remains a significant challenge due to the complexity and multimodal nature of video data. In this paper, we propose \sysname, which, to the best of our knowledge, is the first system to enable VLM-powered video analytics by effectively addressing the aforementioned challenges.








